% ==================== 5 实际测量限制与应对 ====================
\section{实际测量限制与应对}

\noindent\textbf{水的干扰：}红外最棘手，O—H 伸缩（3400 cm$^{-1}$）与 1640 cm$^{-1}$ 弯曲会掩盖 COO$^-$ 信息；应对是样品干燥成膜、ATR 保持干燥并扣背景\cite{piao2026ir}。

\noindent\textbf{样品制备：}KBr 压片易吸潮，还会引入散射与基线倾斜；改用 ATR 或 DRIFTS 可少制样，但压力过大会压坏包覆层。

\noindent\textbf{荧光背景：}油酸杂质或缺陷会产生荧光，可能淹没 144 cm$^{-1}$ 以外的弱峰；785 nm 比 532 nm 荧光低，但信号按 $\lambda^{-4}$ 下降，可用长波长与基线拟合抑制。

\noindent\textbf{信号强度：}拉曼截面极小，COO$^-$ 尤弱，需高数值孔径物镜与多次累加；SERS 虽能增强信号，但基底会改变表面化学\cite{wang2021}。

\noindent\textbf{激光损伤：}功率过高会使油酸脱附、局部升温，甚至诱发锐钛矿→金红石相变或 Ti$^{4+}$→Ti$^{3+}$ 还原\cite{taudul2023}；表面功率宜控制在 1 mW 以下并散焦。

\noindent\textbf{空间分辨率：}常规红外受衍射限制约 1.5--10 $\mu$m，拉曼约 250--500 nm；分辨单颗粒需 AFM-IR/TERS（10--50 nm 与 3--5 nm），但设备复杂且可能扰动包覆层\cite{kurouski2020}；定量还需纯相标样或 PLS 校准\cite{sacco2023}。
